\documentclass[../../main.tex]{subfiles}
\begin{document}

{\bf [解]}
自然数$n$に対して与えられた数列
\begin{align}
 a_n=\displaystyle\int_1^e(\log x)^n\,dx
\end{align}
について考える．

\paragraph{(1)}

$n\ge3$とする．
\begin{align}
 \left(\log x\right)^{n} =  \left(\log x\right)^{n-1}\cdot \log x
\end{align}
とみなして部分積分を用いると
\begin{align}
 a_{n}
  &=\int_1^e(\log x)^{n-1} \log x\, dx \\
 &=\left[x(\log x-1)\left(\log x\right)^{n-1}\right]_{1}^{e} - (n-1)\int_1^e\left\{\left(\log x\right)^{n-1}-\left(\log x\right)^{n-2}\right\} \\
 &= (n-1)\left(A_{n-2}-A_{n-1}\right)
\end{align}
となり，題意は示された．

\paragraph{(2)}

$[1,e]$では$0\le\log x\le1$だから
\begin{align}
 0 \le \left(\log x\right)^{n+1} \le \left(\log x\right)^{n} \label{eq:1}
\end{align}
だから，同区間で積分して
\begin{align}
 0 < a_{n+1} < a_{n} \label{eq:2}
\end{align}
である．ただし，\cref{eq:1}の等号は全ての$x$で成立するわけではないので積分した\cref{eq:2}では等号は除去した．


\paragraph{(3)}

$n\ge2$とする．(1)で$n\to 2n$と置き換えると，
\begin{align}
a_{2n}=(2n-1)(a_{2n-2}-a_{2n-1}) \label{eq:3}
\end{align}
である．ここで(2)から右辺は
\begin{align}
 a_{2n-2}-a_{2n-1}<a_{2n-2}-a_{2n}
\end{align}
と上から評価できるから\cref{eq:3}に代入して
\begin{align}
 a_{2n} &< (2n-1)\left(a_{2n-2}-a_{2n}\right) \\
 2n a_{2n} &< (2n-1)a_{2n-2} \\
\therefore
 a_{2n} &< \dfrac{2n-1}{2n}a_{2n-2}
\end{align}
だから$n,n-1,\cdots,2$について繰り返し用いて
\begin{align}
  a_{2n}<\dfrac{2n-1}{2n}\cdot\dfrac{2n-3}{2(n-1)}\cdots\dfrac{3}{4}\,a_2 \label{eq:4}
\end{align}
である．

また，$a_2$は部分積分によって
\begin{align}
 a_2
  &=\bigl[x(\log x)^2-2x\log x+2x\bigr]_1^e \\
  &=e-2
\end{align}
と求まるから，\cref{eq:4}に代入して
\begin{align}
a_{2n}<\dfrac{(2n-1)(2n-3)\cdots3}{(2n)(2n-2)\cdots4}(e-2)=\dfrac{3\cdot5\cdots(2n-1)}{4\cdot6\cdots(2n)}(e-2)
\end{align}
である．よって題意は示された．

{\bf [解説]}

% TODO :: log^n xの被積分関数について

\newpage

\end{document}
